\documentclass[aspectratio=169,10pt,english]{beamer}

\input{../../../_preambulo_beamer.tex}
\input{../../../_comandos_beamer.tex}

\selectlanguage{english}

\title{MA1001B - Statistical Modeling for Decision Making}
\subtitle{Lesson 39: Two Proportions in A/B Tests}
\author{}
\institute{Tecnol\'ogico de Monterrey}
\date{}

\begin{document}

\begin{frame}[plain]
  \vspace{1.2cm}
  {\Large\bfseries MA1001B - Statistical Modeling for Decision Making\par}
  \vspace{0.7cm}
  {\large Lesson 39: Two Proportions in A/B Tests\par}
  \vspace{0.7cm}
  {\color{TecRojo}\rule{\textwidth}{1.2pt}}
  \vspace{0.8cm}

  MA1001B Faculty\\[0.3cm]
  Term\\[0.3cm]
  Tecnol\'ogico de Monterrey
\end{frame}

\begin{frame}{The Two-Version Conversion Question}
  \begin{alertblock}{Guiding question}
    How do we compare conversion rates from two independent A/B arms, and
    what goes wrong if we peek at $p$ mid-experiment?
  \end{alertblock}

  \vspace{0.6em}
  By the end of this lesson, you can:
  \begin{itemize}
    \item compute $\hat{p}_A$, $\hat{p}_B$, and the pooled $\hat{p}$;
    \item form the $z$ statistic for $p_A-p_B$;
    \item decide $H_0:p_A=p_B$ from a two-sided p-value;
    \item explain why peeking inflates type I error.
  \end{itemize}

  \vfill
  \scriptsize All visitors and conversion flags used today are synthetic. No real company data are used.
\end{frame}

\begin{frame}{Two Independent Conversion Arms}
  \small
  A synthetic checkout test records $x_A=72$ conversions in $n_A=400$
  visitors and $x_B=61$ conversions in $n_B=410$ visitors.

  \[
    \hat{p}_A=\frac{72}{400}=0.18,\qquad
    \hat{p}_B=\frac{61}{410}=0.148780,\qquad
    \hat{p}=\frac{133}{810}=0.164198.
  \]

  \begin{exampleblock}{Two-sided hypotheses}
    $H_0:p_A=p_B$ versus $H_1:p_A\neq p_B$, with $\alpha=0.05$. The arms are
    independent samples, not paired visitors.
  \end{exampleblock}
\end{frame}

\begin{frame}[fragile]{Step 1: Sample Rates}
  \begin{columns}[T]
    \begin{column}{0.42\textwidth}
      \small
      \begin{lstlisting}
phat_a = 72 / 400
phat_b = 61 / 410
p_pool = (72 + 61) / 810
      \end{lstlisting}

      \begin{block}{Verified output}
        $\hat{p}_A=0.180000$\\
        $\hat{p}_B=0.148780$\\
        $p_{\mathrm{pool}}=0.164198$
      \end{block}
    \end{column}
    \begin{column}{0.58\textwidth}
      \centering
      \includegraphics[width=\textwidth]{../figures/two_proportion_01_sample_rates.png}
      \vspace{0.2em}
      \scriptsize Observed conversion rates from Step 1
    \end{column}
  \end{columns}
\end{frame}

\begin{frame}{Pooled $z$ for $p_A-p_B$}
  \small
  \[
    z=\frac{\hat{p}_A-\hat{p}_B}{\sqrt{\hat{p}(1-\hat{p})(1/n_A+1/n_B)}}
    =\frac{0.031220}{0.026035}=1.199141.
  \]
  The two-sided p-value is $0.230473$. Because $|z|<z^*=1.959964$ and
  $p>\alpha$, the test does not reject $H_0:p_A=p_B$.
\end{frame}

\begin{frame}[fragile]{Step 2: Pooled $z$ Decision}
  \begin{columns}[T]
    \begin{column}{0.40\textwidth}
      \small
      \begin{lstlisting}
se = sqrt(p*(1-p)*(1/n_a + 1/n_b))
z = (phat_a - phat_b) / se
p = 2 * norm.sf(abs(z))
      \end{lstlisting}

      \begin{block}{Verified output}
        $SE=0.026035$, $z=1.199141$\\
        $p=0.230473$\\
        Decision: do not reject $H_0$
      \end{block}
    \end{column}
    \begin{column}{0.60\textwidth}
      \centering
      \includegraphics[width=\textwidth]{../figures/two_proportion_02_pooled_z.png}
      \vspace{0.2em}
      \scriptsize Two-sided $z$ tails from Step 2
    \end{column}
  \end{columns}
\end{frame}

\begin{frame}{Step 3: Peeking Inflates Type I Error}
  \begin{columns}[T]
    \begin{column}{0.42\textwidth}
      \small
      Seed-42 simulation of $8000$ null A/B tests, each with a shared
      conversion probability $0.16$.

      \vspace{0.4em}
      One final look: type I rate $=0.047250$.

      \vspace{0.3em}
      Four peeks (25, 50, 75, 100 percent): type I rate $=0.124250$.

      \vspace{0.5em}
      \textbf{Limit:} stopping at the first $p<0.05$ is not an
      $\alpha=0.05$ test.
    \end{column}
    \begin{column}{0.58\textwidth}
      \centering
      \includegraphics[width=\textwidth]{../figures/two_proportion_03_peeking_limit.png}
      \vspace{0.2em}
      \scriptsize Peeking type I inflation from Step 3
    \end{column}
  \end{columns}
\end{frame}

\begin{frame}{Concept Check}
  \begin{enumerate}
    \item Compute $\hat{p}_A$ from $x_A=72$ and $n_A=400$.
    \item Why does the test SE use a pooled $\hat{p}$ rather than two
      separate Wald SEs?
    \item If $p=0.230473$ and $\alpha=0.05$, what is the A/B decision?
    \item Why does four-look peeking raise the false-positive rate to
      $0.124250$?
  \end{enumerate}

  \vspace{0.6em}
  \begin{block}{Connection to Lesson 40}
    The session can continue directly into a chi-square goodness-of-fit test
    for more than two categories.
  \end{block}

  \vfill
  \scriptsize Source: OpenStax, \textit{Introductory Business Statistics 2e},
  Chapter 10, Section 10.4. CC BY-NC-SA.
\end{frame}

\appendix



\end{document}
